% Book-local P03.1 scalar-fiber maximum proof,
% 19 September 2026.  No novelty or priority claim is made.

\begin{proof}[\Cref{lem:p031-perfect-pullback-usc}]
Write
\[
 a^*(x)=\inf_{U\ni x}\sup_U a,
 \qquad
 b(y)=\inf_{N\ni y}\sup_N(a\circ q),
\]
where the infima may be taken over open neighborhoods.  Fix $x\in X$.
For $y\in q^{-1}(x)$ and every neighborhood $U$ of $x$, the set
$q^{-1}(U)$ is a neighborhood of $y$, and surjectivity gives
\[
 \sup_{q^{-1}(U)}(a\circ q)=\sup_Ua.
\]
Taking first the infimum over the neighborhoods of $y$, and then the
infimum over $U$, yields
\begin{equation}
\label{eq:p031-fiber-easy-inequality}
                         b(y)\leq a^*(x).
\end{equation}

The function $b$ is upper semicontinuous.  Its restriction to the compact
fiber $F_x:=q^{-1}(x)$ therefore attains a maximum; denote that maximum by
$M$.  We prove $a^*(x)\leq M$.  Let $r\in\mathbf R$ satisfy $r>M$.
For each $y\in F_x$, the defining neighborhood formula for $b(y)<r$
gives an open neighborhood $N_y$ such that
\[
                         a\circ q<r\quad\text{on }N_y.
\]
Compactness supplies $y_1,\ldots,y_m\in F_x$ for which
$
 N:=N_{y_1}\cup\cdots\cup N_{y_m}
$
contains $F_x$.  Since $q$ is closed,
\[
                         U:=X\setminus q(Y\setminus N)
\]
is an open neighborhood of $x$, and $q^{-1}(U)\subseteq N$.  If
$x'\in U$, surjectivity lets us choose $y'\in q^{-1}(x')$; then
$y'\in N$, so $a(x')=(a\circ q)(y')<r$.  Thus
$
 \sup_Ua\leq r
$
and $a^*(x)\leq r$.  Letting $r\downarrow M$ proves the desired
inequality when $M\in\mathbf R$.  If $M=-\infty$, the same argument
works for every real $r$ and gives $a^*(x)=-\infty$.

Together with \eqref{eq:p031-fiber-easy-inequality}, this proves
\eqref{eq:p031-perfect-pullback-usc}; substituting $x=q(y)$ gives
\eqref{eq:p031-fiber-maximum-pullback}.
\end{proof}

\begin{proof}[\Cref{prop:p031-open-map-characterization}]
If $q$ is open, the neighborhood images $q(N)$, with $N$ ranging over
the neighborhoods of $y$, are cofinal among the neighborhoods of
$q(y)$.  The neighborhood-filter formula therefore gives
$
 (a\circ q)^*(y)=a^*(q(y))
$
for every locally bounded-above $a$; this is
\Cref{lem:p031-open-pullback-usc}.  Hence (1) implies (2), and (2)
implies (3).

Conversely, suppose that $q$ is not open.  There are an open set
$N\subseteq Y$ and a point $y\in N$ such that $q(N)$ is not a
neighborhood of $x:=q(y)$.  Define
\[
                  a:=\mathbf 1_{X\setminus q(N)}.
\]
Every neighborhood of $x$ meets $X\setminus q(N)$, so $a^*(x)=1$.
On the other hand, $a\circ q=0$ throughout the neighborhood $N$ of
$y$, and hence
\[
                  (a\circ q)^*(y)=0<a^*(q(y)).
\]
Thus (3) fails.  This proves that (3) implies (1) and completes the
equivalence.
\end{proof}

\begin{proof}[\Cref{cor:p031-scalar-fiber-maximum}]
Scalar projection is continuous, closed, and surjective by
\cite[\S2.9]{GR25Classical}.  Its fiber over $x$ is canonically
\[
       \mathcal M\bigl(\mathcal H(x)\widehat\otimes_KK'\bigr),
\]
as recalled in \Cref{sec:berk-functoriality}; it is compact by
\Cref{prop:berk-compact-spectrum}.  The hypotheses of
\Cref{lem:p031-perfect-pullback-usc} therefore hold and give
\eqref{eq:p031-scalar-fiber-maximum}.

If $W\subseteq X$ is ordinary open, the restricted map
$
 \pi^{-1}(W)\to W
$
is still continuous and surjective and has the same fibers.  It is closed:
if $F=G\cap\pi^{-1}(W)$ is relatively closed, with $G$ closed in
$X_{K'}$, then
$
 \pi(F)=\pi(G)\cap W
$
is relatively closed in $W$.  The same lemma applies.
\end{proof}

\begin{proof}[\Cref{cor:p031-fiber-constancy-upgrade}]
Apply \Cref{cor:p031-scalar-fiber-maximum} to $a=s$.  Since
$u=s^*$, the result is exactly
\eqref{eq:p031-perron-fiber-maximum}.  If $v_{K'}$ is constant on a
fiber, its common value is its maximum there, so
\eqref{eq:p031-fiber-constancy} holds.  Conversely, if that identity
holds, its left side is constant on each scalar fiber; hence $v_{K'}$ is
fiberwise constant.  This proves the equivalence.

If the condition holds after every analytic extension, the exact
compatibility hypothesis of
\Cref{cor:p031-scalar-compatibility-upgrade} holds.  That corollary gives
classical plurisubharmonicity of $u$, and its obstacle specialization gives
the final assertion.
\end{proof}

\begin{proof}[
\Cref{prop:p031-type-four-topological-defect}]
Poineau--Turchetti take an algebraically closed complete field $K$ which
is not spherically complete, a maximally complete algebraically closed
extension $K'/K$, and a type--$4$ point $x$ of the Berkovich affine
line.  By \cite[Lemma~5.6.2]{PT21SchottkyI}, the fiber over $x$ is a
positive-radius closed disc $D^+(\gamma,r(x))$.  The open disc
$D^-(\gamma,r(x))$ lies in that fiber and therefore has image $\{x\}$,
which is not a neighborhood of $x$; this is the non-openness recorded in
\cite[Remark~5.6.4]{PT21SchottkyI}.

Apply the contrapositive construction in the proof of
\Cref{prop:p031-open-map-characterization} to this open disc and a point
$y$ in it.  It gives a bounded $\{0,1\}$-valued function $a$ such that
\[
                 (a\circ\pi_{K'/K})^*(y)=0<1=a^*(x).
\]
Finally, the affine-line scalar projection is proper and surjective by
\cite[Proposition~5.6.1]{PT21SchottkyI}, so
\Cref{lem:p031-perfect-pullback-usc} identifies $a^*(x)$ with the
fiber maximum and proves \eqref{eq:p031-type-four-strict-defect}.  The
construction imposes no psh property on $a$, which proves the stated
scope boundary.
\end{proof}
